%
% 27-pol.tex -- Der Fall 1/(x-xi)^k\vartheta
%
% (c) 2026 Prof Dr Andreas Müller
%
\subsection{Integrale von $1/(x-\xi)^k\vartheta$
\label{buch:intalg:wurzel:pol:subsection:ktheta}}
Es muss noch gezeigt werden, dass mit den bereits entwickelten
Methoden auch das zweite Integral
in \eqref{buch:intalg:wurzel:eqn:grundintegrale},
\[
\int
\frac{1}{(x-\xi)^k\vartheta}
\,dx,
\]
behandelt werden kann.

Die Substitution
\begin{align*}
x &= \xi + \frac{1}{t} \qquad\text{mit}\qquad dx = -\frac{1}{t^2}\,dt
\intertext{verwandelt den Radikanden in}
ax^2+bx+c
&=
\frac{(a\xi^2+b\xi+c)t^2 + (2a\xi+b)\xi + a}{t^2}
\intertext{und damit den Integranden in}
\frac{1}{(x-\xi)^k\vartheta}
&=
\frac{t^k\cdot t}{
\!\sqrt{
(a\xi^2+b\xi+c)t^2 + (2a\xi+b)t + a
}
}.
\intertext{Das Integral wird damit zu}
\int\frac{1}{(x-\xi)^k\vartheta}\vartheta\,dx
&=
-
\int
\frac{t^k\cdot t}{
\!\sqrt{
(a\xi^2+b\xi+c)t^2 + (2a\xi+b)t + a
}}
\frac{1}{t^2}\,dt
\intertext{Schreibt man $A=a\xi^2+b\xi+c$, $B=2a\xi+b$ und $C=a$, wird
das Integral zu}
&=
-\int\frac{t^{k-1}}{\!\sqrt{At^2+Bt^2+C}}\,dt.
\intertext{Dies ist wieder ein Integral}
&=
-\int \frac{t^{k-1}}{\tilde{\vartheta}}\,dt
\end{align*}
der Form
$\text{Polynom}/\tilde{\vartheta}$, aber diesmal mit
\[
\tilde{\vartheta}
=
\!\sqrt{At^2+Bt+C}.
\]
Diese Art von Integral wurde in
Satz~\ref{buch:intalg:wurzel:satz:polynomtheta}
bereits auf das Integral von $1/\tilde{\vartheta}$
zurückgeführt.

